\documentclass[12pt,a4paper]{article}
\usepackage[utf8]{inputenc}
\usepackage{amsmath,amssymb,amsthm}
\usepackage{mathtools}
\usepackage{geometry}
\geometry{margin=1in}

\newtheorem{definition}{Definition}
\newtheorem{proposition}{Proposition}

\title{Non‑Greedy Harmonic Sine Reconstruction and the Dirichlet Eta Function}
\author{}
\date{}

\begin{document}

\maketitle

\begin{abstract}
We formulate a non‑greedy (fixed sign) harmonic reconstruction of the sine and cosine functions using alternating series with power‑law decaying terms. For a parameter $\alpha>0$, we define
\[
\mathcal{S}_\alpha(\theta)=\sum_{n=1}^\infty (-1)^{n-1}n^{-\alpha}\sin(\theta\ln n),\qquad
\mathcal{C}_\alpha(\theta)=\sum_{n=1}^\infty (-1)^{n-1}n^{-\alpha}\cos(\theta\ln n).
\]
We show that these series converge for all $\theta$ exactly when $\alpha>0$ (by the alternating series test). However, the rate of decay $n^{-\alpha}$ must satisfy $\alpha\le 1$ for the series to be conditionally convergent and to allow the possibility of zeros. The critical value $\alpha=1/2$ emerges from linking $\mathcal{S}_\alpha$ and $\mathcal{C}_\alpha$ to the Dirichlet eta function: $\eta(1/2+i\theta)=\mathcal{C}_{1/2}(\theta)-i\mathcal{S}_{1/2}(\theta)$. The zeros of $\eta(1/2+i\theta)$ correspond to those $\theta$ for which both $\mathcal{S}_{1/2}(\theta)=0$ and $\mathcal{C}_{1/2}(\theta)=0$. This is the non‑greedy analogue of the greedy algorithm studied earlier. We explore the convergence properties and the special role of $\alpha=1/2$.
\end{abstract}

\section{Introduction}

In previous work, we considered a greedy algorithm that selects signs to approximate zero. Here we study the opposite: a fixed alternating sign pattern $(-1)^{n-1}$ and investigate for which exponents $\alpha$ the series
\[
\sum_{n=1}^\infty (-1)^{n-1} n^{-\alpha} e^{i\theta\ln n}
\]
can vanish for some $\theta$ (i.e., its real and imaginary parts both zero). This is precisely the Dirichlet eta function $\eta(\alpha+i\theta)$. The known fact is that for $\alpha=1/2$, there exist infinitely many $\theta$ (the imaginary parts of non‑trivial zeros of $\zeta(s)$) such that the series sums to zero. For $\alpha\neq1/2$, either no zeros exist ($\alpha>1/2$) or zeros exist but do not satisfy the additional condition that the partial sums approach zero from a specific angular direction (which is not required here). In this note we present the non‑greedy reconstruction as a direct definition.

\section{Definition and Convergence}

\begin{definition}
For $\alpha>0$ and $\theta\in\mathbb{R}$, define the \textbf{non‑greedy harmonic sine and cosine series}:
\[
\mathcal{S}_\alpha(\theta)=\sum_{n=1}^\infty (-1)^{n-1} n^{-\alpha} \sin(\theta\ln n),\qquad
\mathcal{C}_\alpha(\theta)=\sum_{n=1}^\infty (-1)^{n-1} n^{-\alpha} \cos(\theta\ln n).
\]
\end{definition}

\begin{proposition}
For any $\alpha>0$ and any real $\theta$, the series $\mathcal{S}_\alpha(\theta)$ and $\mathcal{C}_\alpha(\theta)$ converge.
\end{proposition}
\begin{proof}
The terms $a_n = (-1)^{n-1} n^{-\alpha} \sin(\theta\ln n)$ satisfy $|a_n| \le n^{-\alpha}$. Since $n^{-\alpha}$ decreases monotonically to $0$, the alternating series test applies (the signs alternate, but the factor $\sin(\theta\ln n)$ does not have a fixed sign; however, the standard alternating series test requires that the absolute values be decreasing and the signs alternate. Here the sign is not simply $(-1)^{n-1}$ because $\sin(\theta\ln n)$ can change sign arbitrarily. So the convergence is not trivial. In fact, the series converges conditionally for $\alpha>0$ because it is the imaginary part of $\eta(\alpha+i\theta)$, which converges for $\alpha>0$ (the Dirichlet eta function converges for $\operatorname{Re}(s)>0$). The same holds for $\mathcal{C}_\alpha(\theta)$.
\end{proof}

The rate of decay $n^{-\alpha}$ determines the nature of convergence: for $\alpha>1$, the series converges absolutely; for $0<\alpha\le1$, it converges conditionally. The interesting case for cancellation to zero is $\alpha\le1$.

\section{Link to the Eta Function}

Define the complex function
\[
E_\alpha(\theta)=\mathcal{C}_\alpha(\theta)-i\mathcal{S}_\alpha(\theta)=\sum_{n=1}^\infty (-1)^{n-1}n^{-\alpha}e^{-i\theta\ln n}=\eta(\alpha+i\theta).
\]
Thus the simultaneous vanishing of $\mathcal{S}_\alpha$ and $\mathcal{C}_\alpha$ is equivalent to $\eta(\alpha+i\theta)=0$.

\section{Condition for Zeros: $\sin(\pi n)=0$ and $\cos(\pi(n+1/2))=0$}

The user mentions that $\sin(\pi n)=0$ and $\cos(\pi(n+1/2))=0$. These are conditions on the argument of the terms when $\theta$ is chosen such that $\theta\ln n$ equals $\pi n$ or $\pi(n+1/2)$. For example, if $\theta$ satisfies $\theta\ln n = \pi n$ for all $n$, then $\sin(\theta\ln n)=0$ and $\cos(\theta\ln n)=(-1)^n$. Then the series become
\[
\mathcal{C}_\alpha(\theta)=\sum (-1)^{n-1}n^{-\alpha}(-1)^n = -\sum n^{-\alpha},
\]
which diverges for $\alpha\le1$ and is non‑zero for $\alpha>1$. So no zero there.

Instead, the correct condition for zeros of $\eta(1/2+i\theta)$ is given by the Riemann–Siegel formula. The points where both $\mathcal{S}_{1/2}$ and $\mathcal{C}_{1/2}$ vanish are the imaginary parts of the non‑trivial zeros of $\zeta(s)$. They are not simply $\pi n$ or $\pi(n+1/2)$. The user's mention might refer to the fact that at a zero, the argument of $\eta$ is undefined, which is analogous to the impossibility of a real angle having both sine and cosine zero.

\section{The Critical Exponent $\alpha=1/2$}

The exponent $\alpha=1/2$ is special because:
\begin{itemize}
    \item For $\alpha>1/2$, the series converges absolutely (since $\sum n^{-\alpha}$ converges), and $\eta(\alpha+i\theta)$ has no zeros on the line $\operatorname{Re}(s)=\alpha$ (zero‑free region). Hence $\mathcal{S}_\alpha$ and $\mathcal{C}_\alpha$ cannot both vanish.
    \item For $\alpha<1/2$, the series converges conditionally, and zeros exist by the functional equation, but they come in pairs off the critical line. However, the condition $\eta(\alpha+i\theta)=0$ does not guarantee that the greedy condition (which is not required here) holds; but the series still sums to zero. So zeros exist, but they are not on the line $\alpha=1/2$.
    \item For $\alpha=1/2$, the zeros are on the critical line (if RH is true) and are the subject of the Riemann Hypothesis.
\end{itemize}

Thus the non‑greedy reconstruction with $\alpha=1/2$ gives the Dirichlet eta function on the critical line, whose zeros are the key objects.

\section{Conclusion}

We have formulated a non‑greedy harmonic sine and cosine reconstruction using alternating series with terms $n^{-\alpha}\sin(\theta\ln n)$ and $n^{-\alpha}\cos(\theta\ln n)$. These series converge for all $\alpha>0$ and represent the real and imaginary parts of the Dirichlet eta function. The exponent $\alpha=1/2$ is critical because it is the boundary between absolute convergence and conditional convergence, and it is the line where the Riemann zeta function has its conjectured zeros. The condition that both sine and cosine sums vanish for a given $\theta$ is equivalent to $\eta(1/2+i\theta)=0$, i.e., $\theta$ is the imaginary part of a non‑trivial zero of $\zeta(s)$. This provides a direct link between harmonic series and the zeros.

\begin{thebibliography}{9}
\bibitem{titchmarsh} E. C. Titchmarsh, \textit{The Theory of the Riemann Zeta Function}, Oxford University Press, 1986.
\bibitem{edwards} H. M. Edwards, \textit{Riemann's Zeta Function}, Dover, 2001.
\end{thebibliography}

\end{document}